% SOLUCIÓN Ejercicio 2
\documentclass[12pt,a4paper]{article}
\usepackage[utf8]{inputenc}
\usepackage[T1]{fontenc}
\usepackage[spanish]{babel}
\usepackage{amsmath,amssymb}
\usepackage{booktabs}
\usepackage{siunitx}
\sisetup{locale=FR, output-decimal-marker={,}}
\usepackage{circuitikz}
\usepackage[hidelinks]{hyperref}
\usepackage{pgfplots}\pgfplotsset{compat=1.18}

\begin{document}
\section{Datos experimentales}
La Tabla~\ref{tab:carga} reúne las mediciones.
\begin{table}[htbp]
  \centering
  \caption{Carga de un capacitor}
  \label{tab:carga}
  \begin{tabular}{S[table-format=2.0] S[table-format=1.1]}
    \toprule
    {Tiempo (\si{\milli\second})} & {Voltaje (\si{\volt})} \\
    \midrule
     0 & 0,0 \\
    10 & 2,1 \\
    20 & 3,5 \\
    30 & 4,3 \\
    40 & 4,7 \\
    50 & 4,9 \\
    \bottomrule
  \end{tabular}
\end{table}

\section{Modelo}
\begin{equation}
  V_c(t) = V_{in}\left(1 - e^{-t/\tau}\right), \qquad V_{in}=\SI{5}{\volt}
  \label{eq:carga}
\end{equation}

\section{Estimación de $\tau$}
En $t=\tau$, la ecuación~\eqref{eq:carga} da $V_c = 0{,}632\,V_{in} = \SI{3,16}{\volt}$.
Interpolando linealmente entre 10 y \SI{20}{\milli\second} (Tabla~\ref{tab:carga}):
\[
  \tau \approx 10 + 10\cdot\frac{3{,}16-2{,}1}{3{,}5-2{,}1}
       \approx \SI{17,6}{\milli\second}.
\]

\begin{figure}[htbp]
  \centering
  \begin{tikzpicture}
    \begin{axis}[width=0.75\linewidth,height=5.5cm,xlabel={$t$ (ms)},ylabel={$V_c$ (V)},grid=both,xmin=0,ymin=0,legend pos=south east]
      \addplot[thick,domain=0:50,samples=60]{5*(1-exp(-x/17.6))};
      \addplot[only marks,mark=*] coordinates {(0,0)(10,2.1)(20,3.5)(30,4.3)(40,4.7)(50,4.9)};
      \legend{Modelo con $\tau=17{,}6$ ms, Datos}
    \end{axis}
  \end{tikzpicture}
  \caption{Datos y modelo de carga}\label{fig:carga}
\end{figure}
\end{document}
